\documentclass[11pt]{article}
\usepackage[margin=1in]{geometry}
\usepackage{amsmath,amssymb,amsthm}
\usepackage[hidelinks]{hyperref}
\setlength{\emergencystretch}{2em}
\newtheorem{theorem}{Theorem}
\newtheorem{proposition}[theorem]{Proposition}
\newtheorem{corollary}[theorem]{Corollary}
\newcommand{\mr}{\operatorname{mr}}
\newcommand{\Z}{\mathbb Z}
\newcommand{\F}{\mathbb F}
\newcommand{\rk}{\operatorname{rank}}
\newcommand{\OO}{\mathcal O}
\title{Quadratic coefficients and rigidity of icosahedral lifts}
\author{Internal research working proof}
\date{16 November 2026}
\begin{document}
\maketitle

The retained first nonconstant $A_5$ core colourings of twelve current
targets extend to $H_3$ reflection quotients. They therefore provide
rank-three bases outside the symmetric-transposition route already
excluded. This note gives an exact all-characteristic lift test over
$\OO=\Z[\tau]$, $\tau^2=\tau+1$, and proves that every restricted affine
lift of these fixed bases factors through the same finite group of
order 120. Their lift bound is exactly three in every field; none
gives a new sharp four-bound. Other base colourings are not excluded.
The full Meridional Rank Conjecture remains open.

The uniform arguments are written working proofs. Finite arithmetic
and group certificates have separate replay, with shared routines
identified below. No formal verification, external specialist review
or new-priority claim is made. The use of finite Coxeter quotients is
established, including $H_3$ in the searches of
Blair--Kjuchukova--Morrison \cite{bkm}; the degree-two core-group
framework is due to Kelly and Wada \cite{core}.

\section{From a double-cover map to an icosahedral base}

Write $G=G_K$ and fix a meridian $\mu_0$. The orbifold quotient
$G/\langle\!\langle\mu_0^2\rangle\!\rangle$ is the semidirect product
$H\rtimes\langle\mu_0\rangle$, where
$H=\pi_1(\Sigma_2K)$ and $\mu_0$ has order two. With
$x_i=\mu_i\mu_0$, conjugation by $\mu_0$ sends $x_i$ to $x_i^{-1}$.
The diagram relations on these elements are
\[
                  x_b=x_o x_a^{-1}x_o,\qquad x_0=1.
\]
These are the normalized core relations \cite{core}.

Suppose a retained homomorphism $f:H\to A_5$ has arc images $a_i$.
If $c\in A_5$ is an involution satisfying
$c a_i c^{-1}=a_i^{-1}$ for every arc, set
\[
         \mu_i\longmapsto (a_i c,-1)\quad\hbox{in }A_5\times C_2.
\]
The displayed relations verify this homomorphism directly. In the
twelve cases below, $f$ is onto, $c$ is a double transposition, and
each $a_i c$ is a double transposition. The resulting quotient is onto
the finite group $W=A_5\times C_2$, with meridians in its $H_3$
reflection class. The representation used below identifies this
finite group exactly; no approximate polyhedron is used.

The inversion implementer is found by checking all 120 permutations
of five letters on the existing first map, not by searching for new
core colourings. In every applicable case it is unique and even.
Every original crossing relation is then checked again. The first-map
choice is unchanged after a negative result.

\section{An exact integral reflection model}

Put $\tau=(1+\sqrt5)/2$ and
\[
 A=\begin{pmatrix}2&-\tau&0\\-\tau&2&-1\\0&-1&2\end{pmatrix},
 \qquad S_i=I-e_i A_{i,*}.
\]
Thus $S_i$ fixes the hyperplane $A_{i,*}x=0$ and sends $e_i$ to
$-e_i$. Its difference from the identity has rank one after every
homomorphism $\OO\to F$: each row $A_{i,*}$ contains a unit, and
$\tau(\tau-1)=1$. The roots $g e_i$, $g\in W$, likewise stay nonzero
in every residue field, since $g$ is invertible over $\OO$.

The producer enumerates 120 distinct $3\times3$ matrices over $\OO$,
labelled by $(a,\epsilon)\in A_5\times C_2$. It checks every generator
edge in their Cayley graph and obtains fifteen reflections. The
independent replay instead represents multiplication by $\tau$ as
\[
                  T=\begin{pmatrix}0&1\\1&1\end{pmatrix}
\]
and uses ordinary $6\times6$ integer matrices. It verifies group
closure, inverse matrices and every reflection/root transport.

For a fixed meridian representation $\rho:G\to W$, choose transported
primitive roots $u_i$ for its reflection matrices $R_i$. Restricted
affine lifts take
\[
                  \mu_i\longmapsto (a_i u_i,R_i),\qquad a_i\in F.
\]
They form a vector space $Z_\rho(F)$. The conjugation- and
inversion-invariant root-line condition, followed by evaluation on
arbitrary meridional generators, proves
\[
                         \mr(K)\ge\dim_F Z_\rho(F),
\]
as in the companion all-characteristic proof. These images are
involutions, so their affine relator matrix $M_\rho$ has coefficients
in $\OO$ and one scalar column per arc. It is computed without a
floating point embedding of $\sqrt5$.

There is also a useful comparison with the unrestricted cohomology
bound in Gao--Gu--Li--Wang, Lemma 6.2 \cite{gglw}. Its full elementary
proof has now been read; this does not audit their entire paper.

\begin{proposition}\label{prop:comparison}
For a representation on $F^n$ with meridian defect rank $d>0$, let
$z=\dim Z^1(G;F^n)$ and let $\ell$ be the dimension of the cocycles
restricted by $f(\mu)\in\operatorname{im}(\rho(\mu)-I)$. Then
\[
 \mr(K)\ge\left\lceil\frac{\ell}{d}\right\rceil
 \ \ge\ 1+\left\lceil\frac{z-n}{d}\right\rceil.
\]
Thus the unrestricted cocycle estimate cannot improve the restricted
root-lift bound on the \emph{same} representation.
\end{proposition}
\begin{proof}
The root condition on one meridian propagates to all its conjugates
by the cocycle identity. Evaluation on any $w$ meridional generators
injects the restricted space into $w$ vector spaces of dimension $d$,
so $\ell\le wd$. Evaluation on one fixed meridian, followed by
projection to $F^n/\operatorname{im}(\rho(\mu)-I)$, has kernel the
restricted space and image of dimension at most $n-d$. Hence
$z-\ell\le n-d$, proving the comparison after taking ceilings.
\end{proof}

In particular, the fixed $H_3$ bases below have $n=3,d=1,\ell=3$,
so $z\le5$. Switching to their full first-cocycle spaces cannot
produce a sharp four-bound. A different coefficient representation
is required for such a cohomological improvement.

\begin{corollary}\label{cor:original}
The original nondegenerate binary affine-symplectic criterion also
follows directly from restricted cocycles: a transitive action on
$V=\F_2^k$, $k>0$, with meridians
$x\mapsto x+(B(u_i,x)+a_i)u_i$ and nondegenerate alternating $B$ forces
$\mr(K)\ge k+1$.
\end{corollary}
\begin{proof}
Transitivity forces the meridian directions to span $V$, since each
generator changes a point only by a multiple of its direction.
Nondegeneracy then gives $V^G=0$ for the linear action. Its principal
cocycles form a $k$-dimensional subspace of the restricted cocycles.
The given affine action supplies one more restricted cocycle, outside
that subspace: a principal affine action has a common fixed point
and cannot be transitive on $V$ with $k>0$. Hence the restricted
dimension is at least $k+1$, and the rank-one bound applies.
\end{proof}

This is an alternative cohomological derivation of the user's
nondegenerate conditional theorem, without a novelty claim. For a
degenerate form the common linear fixed space is the radical, so
this argument alone gives only $k-\dim\operatorname{rad}B+1$.
The stronger degenerate $k+1$ extension still uses the separate
quadratic-refinement proof retained in \path{PROOF.tex}.

\section{Finite exact reduction for every coefficient field}

For a matrix $M\in\operatorname{Mat}_{m\times s}(\OO)$, replace each
coefficient $a+b\tau$ by its multiplication matrix
\[
                    \begin{pmatrix}a&b\\b&a+b\end{pmatrix}.
\]
Write $E(M)$ for the resulting $2m\times2s$ integer matrix.

\begin{theorem}\label{thm:quadratic}
Let $t=\rk_{\mathbb Q(\sqrt5)}M>0$. The integer matrix $E(M)$ has
rank $2t$. Let $\delta_1\mid\cdots\mid\delta_{2t}$ be its nonzero
Smith factors. To determine the nullity of $M$ over every field
equipped with a homomorphism $\OO\to F$, it suffices to check:
\begin{itemize}
\item the generic nullity $s-t$;
\item every prime divisor of $\delta_{2t}$, in every residue embedding;
\item characteristic five, with $\tau\mapsto3$.
\end{itemize}
For a prime $p\ne5$, the residue fields are $\F_p$ at each of the two
roots of $X^2-X-1$, or $\F_{p^2}$ when this polynomial is irreducible.
Rank is unchanged by any further field extension. If $M$ has rank
zero, it is identically zero and has nullity $s$ everywhere.
\end{theorem}
\begin{proof}
The regular representation of the quadratic field doubles the generic
rank, giving $2t$. If $p\ne5$, the algebra
$\OO\otimes\F_p=\F_p[X]/(X^2-X-1)$ is either the product of two
copies of $\F_p$ or the field $\F_{p^2}$. The rank of $E(M)$ modulo
$p$ is respectively the sum of the two component ranks or twice the
rank over $\F_{p^2}$.

All $(t+1)$-minors of $M$ vanish in $\OO$, so every component rank is
at most $t$. If $p\nmid\delta_{2t}$, the integer Smith form has rank
$2t$ modulo $p$, forcing every component rank to equal $t$. This
leaves only the stated finite set of primes. At $p=5$, the quadratic
algebra is not reduced, but its unique residue field has $\tau=3$;
checking that scalar matrix directly avoids any semisimplicity
assumption. Characteristic-zero embeddings and further field
extensions preserve the corresponding ranks.
\end{proof}

This is a rank-drop criterion before any field search, not a prime
sweep. The implementation stops without an all-characteristic
certificate if its frozen factor or residue-root cap is exceeded.
On this panel every nonzero Smith factor is a unit, so no factoring
or large-prime search is needed. Characteristic five is nevertheless
checked separately.

Controls distinguish the arithmetic cases. The coefficient $\tau-4$
has norm eleven and vanishes only at one of the two embeddings
modulo eleven. The coefficient $2\tau-1$ vanishes in the residue
field of characteristic five; its expanded kernel has dimension one,
which cannot be blindly halved. Inert characteristic-three and unit
controls are also retained.

\section{When a lift cannot enlarge its finite base}

\begin{proposition}\label{prop:universal}
Restricted root lifts of the finite group $W$ itself form $F^3$ over
every field $F$ receiving $\OO$. Their three coordinates are the
translations $a_i e_i$ on $S_i$.
\end{proposition}
\begin{proof}
Over $\mathbb Q(\sqrt5)$, every triple $(a_1,a_2,a_3)$ is obtained by
conjugating the linear action by a translation $b$ satisfying $Ab=a$.
Indeed
\[
            \det A=2(2-\tau)\ne0,\qquad (2-\tau)(1+\tau)=1.
\]
Every group relator therefore vanishes identically in these three
parameters. Its affine translation is linear in them with coefficients
in $\OO$, so the identity also holds over $\OO$ and after every
specialization to $F$. Thus every triple gives a homomorphism from
the same finite group $W$, including in characteristic two. Three
simple-generator images determine a lift, proving the assertion.
\end{proof}

\begin{theorem}\label{thm:factor}
Suppose $\rho:G\twoheadrightarrow W$ is onto and
$\dim Z_\rho(F)=3$. Every restricted affine lift of this fixed $\rho$
factors through $W$. Its image has order at most 120.
\end{theorem}
\begin{proof}
Pullback along $\rho$ maps the three-dimensional space in
Proposition~\ref{prop:universal} into $Z_\rho(F)$. It is injective:
two homomorphisms from $W$ agreeing after a surjective pullback agree
on $W$. Equality of dimensions makes it onto. This remains valid
when the specialized linear representation of $W$ is not faithful.
\end{proof}

For characteristic different from two, $A$ is invertible and these
universal lifts have a common fixed point. In characteristic two,
the Cartan matrix has rank two, so this last conclusion need not
hold. The factorization conclusion still does.

An exact control over $\F_4$ takes simple translation parameters
$(1,0,0)$. It has no common fixed point, image order 120 and four
orbits on $\F_4^3$ of sizes
\[
                              12,12,20,20.
\]
It is not transitive. More generally none of these finite-base lifts
is transitive on $F^3$: for a finite such field, $|F|=q\ge4$;
$4^3$ does not divide 120, and $q^3>120$ for $q\ge5$. For infinite
$F$, a finite group cannot be transitive on $F^3$ either. Over $\F_4$
the meridian difference has rank two on the underlying binary space;
it is not a rank-one binary transvection satisfying the original
quadratic criterion.

\section{The fixed panel and its certificates}

The twelve fixed bases are obtained from the existing first core maps
on
\[
\begin{array}{rrrr}
13a129&13a135&13a136&13a137\\
13a194&13a196&13a197&13a198\\
13a234&13a242&13a244&13a265.
\end{array}
\]
The other four targets in the unchanged sixteen-case panel have only
the trivial $A_5$ core map. No alternative map or knot was substituted.
The panel originates in the pinned October uncertainty list
\cite{deposit}, with independently retained KnotInfo PD data and
four-seed upper certificates from the preceding audit. An uncertainty
list is not evidence of novel MRC status.

For every one of the twelve bases, $M_\rho$ has thirteen columns.
Its expanded integer matrix has 78 rows, 26 columns and nonzero
Smith factors
\[
                                1^{20}.
\]
The generic root nullity is three; at $p=5$, $\tau=3$, independent
scalar elimination again gives three. Theorem~\ref{thm:quadratic}
therefore proves nullity three in every coefficient field, and
Theorem~\ref{thm:factor} proves finite-base factorization in each.

Explicit certificates make this conclusion reviewable. For each
case the retained matrices $P\in\operatorname{Mat}_{13\times3}(\OO)$
and $L\in\operatorname{Mat}_{3\times13}(\OO)$ satisfy
\[
                              M_\rho P=0,\qquad LP=I_3.
\]
Here $P$ is the pullback of the universal three-parameter lift and
$L$ evaluates words mapping onto the three standard generators of
$W$. Those words also certify that the abstract image is all 120
elements. The verifier checks the universal affine Cayley graph,
every pullback coefficient, both displayed identities and the words,
using ordinary integer regular matrices. In characteristic five it
checks a full independent kernel basis. The $\F_4$ orbit control is
replayed on six binary coordinates without quadratic-field arithmetic.

The main scripts are \path{quadratic_integer_lifts.py},
\path{audit_icosahedral_lifts.py} and
\path{verify_icosahedral_lifts.py}. The frozen plan, group model,
root matrices, Smith operations, field bases, inversion audit and
split factorizations are sealed by
\path{icosahedral_lift_verification.json}. The independent replay
does not use the producer's algebraic-number operations or affine
matrix builder; the PD parser and existing Smith-step checker are
shared. These finite replays do not constitute formal verification
of the uniform theorems.

This closes a specific noncrystallographic affine-lift route on the
fixed first maps. It does not exclude all $H_3$ colourings of these
knots, all rank-one representations, nonlinear deformations or
nonabelian lower bounds. No previously unresolved knot case is
settled. Further work should use a structural reason for a different
base or improve geometric bridge upper certificates, while auditing
known-family coverage and the remaining source proofs.

\begin{thebibliography}{9}
\bibitem{gglw} Gao--Gu--Li--Wang, Meridional ranks of iterated
Whitehead doubles (2026), Lemma 6.2 and Remark 6.3, preprint.
\url{https://arxiv.org/abs/2611.15673v1}.
\bibitem{bkm} Blair--Kjuchukova--Morrison, Coxeter quotients of knot
groups through 16 crossings (2024), Sections 2 and 3.
\url{https://arxiv.org/abs/2208.09032}.
\bibitem{core} Silver--Traldi--Williams, Core groups,
J. Knot Theory Ramifications 33 (2024), 2450013, introduction and
Section 7; Kelly--Wada core-group precedent.
\url{https://arxiv.org/abs/2311.02048v2}.
\bibitem{deposit} Bhattacharjee, The bridge index and the nonorientable
four-genus of knots with at most thirteen crossings, unrefereed
release 1.3.0 (2026). \url{https://doi.org/10.5281/zenodo.23561157}.
\end{thebibliography}
\end{document}
